\section{Regimes, unresolved physics and the route to a design}\label{sec:validity}

\subsection{Scales that determine which equations may be reduced}
Choose a characteristic geometric length $L>0$, slow velocity $U_s$, slow
variation time $T_s$, acoustic velocity amplitude $V_a$, sound speed $c$, and
reference density, viscosity and surface tension $\rho,\mu,\sigma$. These are
representative positive scales, not assigned apparatus values. Let
$\nu=\mu/\rho$ be kinematic viscosity, $D_T$ thermal diffusivity, and
$\Pi_a$ a characteristic acoustic traction. Where gravity is relevant,
$g_{\rm grav}$ is its magnitude, distinct from the boundary source $g$.
Define the following dimensionless ratios and times:
\begin{align}
 \mathrm{Ma}_a&=V_a/c, &
 \mathrm{Re}_s&=\rho U_sL/\mu, &
 \mathcal I_s&=\rho L^2/(\mu T_s),\label{eq:scale-ratios-a}\\
 \mathrm{Oh}&=\mu/\sqrt{\rho\sigma L}, &
 \mathrm{Bo}&=|\Delta\rho|g_{\rm grav}L^2/\sigma, &
 \mathrm{Ca}_a&=\Pi_aL/\sigma,\label{eq:scale-ratios-b}\\
 \tau_\sigma&=\sqrt{\rho L^3/\sigma}, &
 \tau_\mu&=\mu L/\sigma, &
 \tau_\nu&=\rho L^2/\mu,\label{eq:scale-times}\\
 \delta_\nu&=\sqrt{2\nu/\omega}, &
 \delta_T&=\sqrt{2D_T/\omega}.&&\label{eq:boundary-layer-scales}
\end{align}
Here $\mathcal I_s$ measures unsteady inertia relative to viscosity,
$\mathrm{Oh}$ is the Ohnesorge number, $\mathrm{Bo}$ the Bond number, and
$\mathrm{Ca}_a$ an acoustic traction/capillary scale ratio. The three $\tau$
quantities are capillary--inertial, viscous--capillary and momentum-diffusion
times; $\delta_\nu,\delta_T$ are viscous and thermal penetration depths.
Both fluids may supply important scales, so a favorable ratio in one phase
does not justify eliminating the other.

\begin{longtable}{@{}>{\raggedright\arraybackslash}p{34mm}p{111mm}@{}}
\toprule Reduction & Conditions to examine before using it \\\midrule\endhead
Geometric optics & Optical wavelength small relative to relevant curvature
and wavefront scales; intended polarization, dispersion and diffraction
requirements treated separately.\\
Linear acoustics & $\mathrm{Ma}_a\ll1$, small fractional thermodynamic
perturbations, and acoustic particle displacement small relative to local
geometric scales. The eventual mean interface deformation need not be small.\\
Thin acoustic layers & $\delta_\nu,\delta_T$ small compared with relevant wall
curvature, tangential variation and cavity scales for the particular
boundary-layer approximation \cite{bach2018,jorgensen2021}.\\
Cycle averaging & A window satisfying \cref{eq:averaging-separation}; beats,
modulation and optical exposure must also be resolved or justified as negligible.\\
Settled harmonic field & Acoustic ring-down short relative to envelope and
interface changes, with compatible actuator settling.\\
Creeping mean flow & Both $\mathrm{Re}_s\ll1$ and $\mathcal I_s\ll1$ over the
formation and perturbation times being studied. Small convective inertia
alone does not remove capillary oscillations.\\
Hydrostatic target load & Negligible mean flow, compatible tangential balance,
constant density and surface tension, and the declared volume and wetting law.\\
Isothermal optics & Thermally induced shape, wave-transfer and index changes
small compared with the requested tolerances over the full operating time.\\
\bottomrule
\end{longtable}

As an explicitly limited analytic example, consider two inviscid, infinitely
deep fluids separated by a flat horizontal interface, with the denser fluid
below. Let $k_\perp>0$ be an interfacial spatial wavenumber and
$\omega_\Gamma(k_\perp)$ its gravity--capillary angular frequency. Then
\begin{equation}
 \omega_\Gamma^2(k_\perp)
   =\frac{\Delta\rho\,g_{\rm grav}k_\perp+\sigma k_\perp^3}{\rho_-+\rho_+}.
 \label{eq:inviscid-capillary-limit}
\end{equation}
This is a scale check for fast interface compliance, not the eigenfrequency
of an unspecified chamber. Finite depth, wetting, curvature and viscosity
change its response; the viscous problem has a nontrivial dispersion relation
\cite{denner2016}. In particular, making the carrier much faster than the
relevant interfacial modes is one ingredient in neglecting first-order
capillary corrections to acoustic scattering. The actual coupled boundary
problem determines whether that neglect is accurate enough optically.

\subsection{Physical effects that may set the attainable optical error}
No single hierarchy is guaranteed for all materials and arrangements. The
following effects enter the same governing problem when their consequences
are comparable with the requested optical tolerance:
\begin{enumerate}
\item \emph{Streaming and absorption}: steady recirculation changes viscous
tractions, temperature and the acoustic field. It is not equivalent to adding
a fitted normal-pressure offset.
\item \emph{Thermal and composition feedback}: density, viscosity, sound speed,
attenuation, surface tension and optical index can all depend on temperature
or composition. Surface-tension gradients create additional tangential stress.
\item \emph{Fast optical modulation}: interface motion and index modulation in
\cref{eq:fast-optical-error} may affect an exposure even when the mean shape
is exactly Cartesian.
\item \emph{Wetting and rheology}: hysteresis, evaporation, contamination,
non-Newtonian response or curing alter the state variables or the admissible
equilibria. Each requires its own measured or justified closure.
\item \emph{Actuator and boundary loading}: radiation, elastic compliance,
transducer coupling and source calibration affect both attainable loads and
their shape derivatives. A finite array is not an arbitrary pressure boundary.
\item \emph{Disturbances and resolution}: vibration, parameter drift and finite
sensor bandwidth act on the actual dynamic system. Tolerances must account for
their optical effect, not only a deterministic design residual.
\end{enumerate}
These are model-selection questions, not reasons to include every effect
indiscriminately. The decision is made by scales and sensitivity against a
stated optical requirement.

Finite actuator authority particularly limits arbitrary fine spatial detail.
For a nearly flat surface with a small sinusoidal height component of
wavenumber $k_\perp$, the capillary part of the required traction scales as
$\sigma k_\perp^2$ times its height amplitude. Thus fine features require large
curvature loads. Whether the available wave fields can provide those loads is
determined by \cref{eq:quadratic-load}, not by the number of pixels in a target
image. Cartesian surfaces form a restricted optical family, which is a useful
design focus, but their membership in that family does not prove actuatability.

\subsection{A theoretical specification before any numerical experiment}
The general problem is specified in this order:
\begin{enumerate}
\item Supply the cylinder, fluid laws, conserved mass or volume, initial state,
wall response and wetting conditions. Choose a class of regular admissible
surfaces without imposing source or target axisymmetry.
\item Specify the physical source supports, commanded variables, constitutive
source response, amplitude, bandwidth and energy limits. Select justified
wave, mean-flow and thermal models and their averaging regime.
\item For a target surface, impose its geometric and volume compatibility.
Derive the required normal load when the quiescent reduction applies;
otherwise retain the full stationary flow and stress equations.
\item Formulate the continuous source inverse, test necessary load bounds and
derive its state, adjoint and constrained optimality equations. Define the
physical tolerance and regularization independently of numerical resolution.
\item Linearize the actual maintained state, including memory and all admitted
three-dimensional disturbances. Determine the required observations and source
authority for feedback, if the open-loop state is unstable.
\item Specify formation from the supplied initial state and the intended terminal
condition: driven operation, temporary transient use, or a separately modeled
retention mechanism. State uncertainty and the one realizable control policy.
\item Specialize the target and observations to Cartesian optics if desired.
The conjugates select an optical surface; they do not replace the preceding
physical data or feasibility conditions.
\end{enumerate}
This construction supplies a reusable inverse formulation. A finite array,
a numerical basis, and a nonlinear solution algorithm are later choices.
Continuation can explore branches and expose folds; it is not the definition
of the inverse map, a required source-generation procedure, or evidence of
coverage of the target family.

\subsection{Numerical experiments will test stated predictions}
The following tests are consequences of the equations, to be performed when
numerical work resumes. None is reported as executed in this revision.

\begin{longtable}{@{}p{36mm}p{108mm}@{}}
\toprule Prediction or identity & Experiment and what a discrepancy would mean
\\\midrule\endhead
Source work, \cref{eq:power-identity} & Compare integrated source power with
boundary absorption using the same solution and units. A persistent difference
indicates a sign, phasor normalization, boundary, or omitted-power-term error.\\
Source adjoint, \cref{eq:source-green-identity,eq:distributed-gradients} &
Check the Green identity and independent directional derivatives for boundary
and volume sources. Include a velocity-dependent interface objective; a
pressure-only check cannot detect a missing adjoint pressure jump.\\
Shape sensitivity, \cref{eq:wave-sensitivity,eq:curvature-variation} &
Perturb geometry while holding the same actual physical command fixed. Update
wave scattering, normals, measure, wall loading and mean flow. Compare the
complete directional derivative under separate refinements.\\
Angular selection, \cref{sec:sources} & In a rotationally invariant geometry,
excite specified angular wave orders and check the predicted load differences
$m-n$. Test asymmetric geometry separately, where modes should couple.\\
Source authority, \cref{eq:continuous-certificates} & Compare optimized loads
with independent lower bounds on source effort and with relaxation bounds.
A violating solution signals an operator, normalization or constraint error.
Failure of an optimizer alone does not prove infeasibility.\\
Passive dynamics, \cref{eq:passive-energy-decay} & Switch off actuation in a
closed passive limit. Check mass and dissipation. A non-capillary stationary
shape with every sustaining effect removed contradicts the retention condition.\\
Coupled stability, \cref{eq:streaming-linearization} & Perturb the complete
holding state, including non-axisymmetric modes, thermal and wave memory when
retained. Compare growth with the linearized prediction inside its local range.\\
Source-space approximation & Refine an admissible source basis with the physical
model fixed. Separately refine wave, interface, flow and time discretizations.
Agreement is evidence for those instances, not a proof of the compactness
hypothesis or the general approximation proposition.\\
Cartesian specialization & Check the unsquared path, the transmitted Snell
branch, analytic curvature and actual ray errors. Include fast interface and
index modulation if they exceed the optical tolerance.\\
\bottomrule
\end{longtable}

Earlier apparatus-specific calculations remain examples under their declared
approximations. They do not set the source space or prove the propositions in
this manuscript. Before interpreting re-optimization on a finer model, establish
convergence with the same physical drive and account for every material loss
mechanism. Experimental validation subsequently needs calibrated source motion,
material properties, measured interface dynamics and optical response.

\subsection{What this formulation establishes and leaves open}
The source-to-wave-to-load map, conditional reachability bounds, regularized
existence argument, explicit acoustic source adjoint, and constrained dynamic
optimality equations provide a common theoretical entry to inverse shaping.
Cartesian surfaces supply a useful target family inside that problem. No
universal excitation formula follows from four optical parameters alone.

The next theoretical specialization must choose constitutive and wall laws
accurate enough for an intended physical experiment, derive the remaining
thermoviscous and moving-interface adjoint terms, and establish the relevant
regularity and stability bounds. Numerics can test approximations and proposed
algorithms, including their failure cases; they cannot promote assumptions
about unmodeled physics into a theorem. A stable machine covering a stated
Cartesian parameter region ultimately requires a documented feasible region,
source selection rule, formation policy and validation throughout that region.
